\section{Критерии качества}

\begin{frame}{Что значит «хорошая модель»?}
  Прогноз спроса должен помогать магазину уменьшить дефицит и списания.
  \begin{description}
    \item[До запуска] Сравниваем прогноз с известным спросом на отложенных примерах: это офлайн-оценка.
    \item[После запуска] Измеряем дефицит, списания и расходы магазина: это онлайн-оценка.
  \end{description}
  \medskip
  Меньшая ошибка прогноза не обязательно означает меньшие расходы магазина.
\end{frame}

\begin{frame}{Офлайн или онлайн?}
  \begin{numrows}
    1 & Доля верных ответов на отложенных фотографиях.\\
    2 & Число списанных товаров после запуска прогноза спроса.\\
    3 & Средняя ошибка предсказания цены на известных продажах.\\
    4 & Доля пользователей, совершивших покупку в эксперименте.\\
  \end{numrows}
  \medskip
  Какие данные нужны для каждой оценки?
\end{frame}

\begin{frame}{Разбор: офлайн- и онлайн-оценка}
  \begin{numrows}
    1 & Офлайн: фотографии и правильные подписи уже собраны.\\
    2 & Онлайн: наблюдаем списания в работающей системе.\\
    3 & Офлайн: цены продаж известны заранее.\\
    4 & Онлайн: измеряем поведение участников эксперимента.\\
  \end{numrows}
  \medskip
  Оценку людьми можно проводить и до запуска, и после него. «Онлайн» означает оценку работающей системы.
\end{frame}

\begin{frame}{Два вида ошибок классификации}
  Модель ищет сообщения со спамом. Положительный класс --- \textbf{спам}.
  \medskip
  \par
  \begin{tabularx}{\linewidth}{@{}Y Y Y@{}}
    \toprule
    & На самом деле спам & Обычное сообщение\\
    \midrule
    Назвала спамом & Верно найденный спам (TP) & Ложная тревога (FP)\\
    Не назвала спамом & Пропущенный спам (FN) & Верное отрицание (TN)\\
    \bottomrule
  \end{tabularx}\par
  \medskip
  Что хуже --- потерять важное письмо или пропустить рекламу? Ответ зависит от цели сервиса.
\end{frame}

\begin{frame}{Accuracy, precision и recall}
  \par
  \begin{tabularx}{\linewidth}{@{}l Y@{}}
    \toprule
    \textbf{Метрика} & \textbf{Какой вопрос задаёт}\\
    \midrule
    Accuracy & Какая доля всех ответов верна? $(TP+TN)/N$.\\
    Precision & Какая доля сообщений, названных спамом, действительно спам? $TP/(TP+FP)$.\\
    Recall & Какую долю настоящего спама нашли? $TP/(TP+FN)$.\\
    \bottomrule
  \end{tabularx}\par
  \medskip
  $N$ --- число всех сообщений. Если знаменатель равен нулю, правило вычисления нужно оговорить отдельно.
\end{frame}

\begin{frame}{Посчитайте метрики}
  Из 100 сообщений 10 --- спам. Модель отметила 8 сообщений: 6 со спамом и 2 обычных.
  \begin{block}{Задание}
    Найдите TP, FP, FN, TN. Затем вычислите accuracy, precision и recall.
  \end{block}
  \medskip
  Какую accuracy получит модель, которая всегда отвечает «не спам»?
\end{frame}

\begin{frame}{Ответ и ловушка accuracy}
  $TP=6$, $FP=2$, $FN=4$, $TN=88$.
  \begin{itemize}
    \item Accuracy: $(6+88)/100=94\%$.
    \item Precision: $6/8=75\%$.
    \item Recall: $6/10=60\%$.
  \end{itemize}
  \medskip
  «Всегда не спам» даёт accuracy $90\%$, но не находит ни одного сообщения со спамом.
  \smallskip\par
  Для редкого положительного класса оценивайте precision и recall вместе с accuracy.
\end{frame}

\begin{frame}{Метрика должна соответствовать цели}
  \begin{itemize}
    \item Цена квартиры: MAE показывает средний модуль ошибки в денежных единицах.
    \item Число людей в помещении: ошибка подсчёта и пропуски превышения допустимого числа --- разные цели.
    \item Рамка вокруг объекта: важны и обнаружение объекта, и точность положения рамки.
  \end{itemize}
  \medskip
  Также учитывают время ответа, память и возможность объяснить решение.
  \medskip
  Перед обучением зафиксируйте цель, метрику и простое решение для сравнения.
\end{frame}
